\Askhsh{34330}{4}
\begin{ylh}{1}
    \item 3.1. Είδη και στοιχεία τριγώνων
    \item 3.2. 1ο Κριτήριο ισότητας τριγώνων
    \item 3.7. Κύκλος - Μεσοκάθετος - Διχοτόμος
    \item 3.10. Σχέση εξωτερικής και απέναντι γωνίας
    \item 5.9. Μια ιδιότητα του ορθογώνιου τριγώνου.
\end{ylh}
% ===================================================================
%  Αρχείο: 34330.tex (Fragment)
%  Αυτόματη μετατροπή μέσω doc2latex.py
% ===================================================================

ΘΕΜΑ 4

Δίνονται τα ορθογώνια τρίγωνα ΑΒΓ ($\widehat{Α}$ = 90\textsuperscript{0} ) και ΔΒΓ ($\widehat{\Delta}$ = 90\textsuperscript{0} ) με τις κορυφές τους Α και Δ εκατέρωθεν της ΒΓ και το μέσο Μ της ΒΓ.

Να αποδείξετε ότι:

\begin{alist}
\item το τρίγωνο ΑΜΔ είναι ισοσκελές, \monades{9}

\item Α$\widehat{Μ}$Δ = 2Α$\widehat{\Gamma}$Δ, \monades{9}

\item τα Α, Β, Δ και Γ είναι σημεία ενός κύκλου, τον οποίο και να κατασκευάσετε.

\monades{7}


\begin{center}
\begin{tikzpicture}[scale=1]
\tkzDefPoint(0,0){B}
\tkzDefPoint(4,0){C}
\tkzDefMidPoint(B,C) \tkzGetPoint{M}

% The circle has diameter BC.
\tkzDrawCircle(M,B)
% A and D are on the circle, opposite sides of BC.
\tkzDefPointBy[rotation=center M angle 110](C) \tkzGetPoint{A}
\tkzDefPointBy[rotation=center M angle -50](C) \tkzGetPoint{D}

\draw[l3] (A)--(B)--(C)--cycle;
\draw[l3] (D)--(B)--(C)--cycle;
\draw[l3] (A)--(M);
\draw[l3] (D)--(M);
\draw[l3] (A)--(D);
\draw[l3] (A)--(C);
\draw[l3] (D)--(C);

\tkzDrawPoints(A,B,C,D,M)
\tkzLabelPoint[above left](A){$A$}
\tkzLabelPoint[left](B){$B$}
\tkzLabelPoint[right](C){$\Gamma$}
\tkzLabelPoint[below right](D){$\Delta$}
\tkzLabelPoint[below](M){$M$}

\tkzMarkRightAngle(B,A,C)
\tkzMarkRightAngle(B,D,C)
\end{tikzpicture}
\end{center}
\end{alist}
